% from ega4-2-fr.tex
\oldpage[IV-2]{152}
\begin{enumerate}
\item[(i)] On sait déjà que si $f$ et $g$ sont plats, il en est de même de
$u=g\circ f$, et que si $f$ et $g\circ f$ sont plats et $f$ surjectif, $g$ est
plat (2.1.6 et 2.2.13). D'autre part, pour tout $z\in Z$, le morphisme
$f_z=f\otimes\mathbf 1_{k(z)}:u^{-1}(z)\to g^{-1}(z)$ est plat (resp.
fidèlement plat si $f$ l'est) et, pour tout $y\in g^{-1}(z)$,
$f_z^{-1}(y)$ est isomorphe à $f^{-1}(y)$ (I, 3.6.4). Si $M(f)$ et $M(g)$
sont vraies, $M(g\circ f)$ est donc vraie pour les cas où $M$ est la propriété
d'être de Cohen--Macaulay ou $(S_n)$, en vertu de (6.6.1). D'autre part soit
$k'$ une extension finie de $k(z)$; posons
$Y'_z=g^{-1}(z)\otimes_{k(z)}k'$, $X'_z=u^{-1}(z)\otimes_{k(z)}k'$, et
$f'_z=f_z\otimes\mathbf 1_{k'}:X'_z\to Y'_z$. Le morphisme $f'_z$ est plat
(resp. fidèlement plat) et, pour tout $y'\in Y'_z$, la fibre $f_z'^{-1}(y')$
est isomorphe à $f_z^{-1}(y)\otimes_{k(y)}k(y')$, en désignant par $y$ l'image
de $y'$ dans $g^{-1}(z)$ (I, 3.6.4). Lorsque $M$ est la propriété d'être
régulier, $(R_n)$, normal ou réduit, l'hypothèse que $M(f)$ et $M(g)$ sont
vraies entraîne que $Y'_z$ et chacune des fibres $f_z'^{-1}(y')$ possèdent
pour $y'\in Y'_z$ la propriété correspondante parmi les propriétés c), d),
e), f) de (6.6.1); on déduit donc de (6.6.1, (i)) que $X'_z$ possède la même
propriété, donc $M(g\circ f)$ est vraie.
\item[(ii)] Inversement, l'hypothèse que $M(g\circ f)$ et $M(f)$ sont vraies
et que $f$ est surjectif entraîne que $Y'_z$ a la propriété correspondante en
vertu de (6.6.1, (ii)), $f'_z$ étant surjectif pour tout $z$; donc $M(g)$ est
vraie.
\item[(iii)] La première assertion découle aussitôt de (6.8.2). D'autre part,
si $h$ est fidèlement plat et $f'$ plat, $f$ est plat (2.4.1); comme, avec les
notations de (6.8.2), les propriétés $P(f^{-1}(y),x)$ et
$P(f'^{-1}(y'),x')$ sont équivalentes (6.7.8), on voit que $M(f)$ et $M(f')$
sont alors équivalentes.
\end{enumerate}
Enfin, la dernière assertion de la proposition résulte de (1.4.3, (iii)) et
de (2.7.1, (iv)).
\end{proof}

\begin{remarks}[6.8.4]
\label{IV.2.s6.8.4-fr}
\begin{enumerate}
\item[(i)] Si $f$ est fidèlement plat, $g$ plat et si $g\circ f$ est de
coprofondeur $\leq n$, alors $g$ est de coprofondeur $\leq n$, comme il
résulte de (6.6.2).
\item[(ii)] Lorsque, dans (6.8.1), on prend pour $Y$ le spectre d'un corps
$k$, les notions (iv), (v) et (vi) se réduisent à celles définies dans (6.7.5).
Il est clair que ces dernières sont relatives au corps de base $k$. La
définition (6.8.1) conduit alors à dire qu'un préschéma $X$ est «régulier
(resp. $(R_n)$, resp. normal) sur $k$» au lieu de dire qu'il est
«géométriquement régulier (resp. $(R_n)$, resp. normal) relativement à $k$»;
on évitera de confondre cette notion avec la propriété d'être régulier (resp.
$(R_n)$, resp. normal) qui est indépendante de $k$. On peut faire à ce propos
les mêmes remarques que dans (4.5.12).
\end{enumerate}
\end{remarks}

\begin{proposition}[6.8.5]
\label{IV.2.s6.8.5-fr}
Soient $k$ un corps, $X$, $Y$ deux $k$-préschémas localement noethériens,
dont l'un est localement de type fini sur $k$. Pour un $k$-préschéma $Z$,
désignons par $P(Z)$ l'une des propriétés c) à f) de (6.6.1); la propriété
«$P(Z)$ géométrique» est alors définie dans (6.7.6) (resp. (4.6.1)) lorsque
$P(Z)$ est l'une des propriétés c), d), e) (resp. f)) de (6.6.1). Alors:
\begin{enumerate}
\item[(i)] Si $X$ possède la propriété $P$, et $Y$ la propriété $P$
géométrique, $X\times_kY$ possède la propriété $P$.
\item[(ii)] Si $X$ et $Y$ possèdent la propriété $P$ géométrique, il en est
de même de $X\times_kY$.
\end{enumerate}
\end{proposition}

\begin{proof}
Soient en effet $f:X\times_kY\to X$, $g:X\to\operatorname{Spec}(k)$ les
morphismes structuraux, qui sont fidèlement plats (2.2.13). L'hypothèse que
$Y$ possède la propriété $P$ géométrique entraîne en vertu de (6.8.2) que
$M(f)$ est vraie, $M$ étant la propriété de (6.8.3) qui correspond à $P$; sous
l'hypothèse (ii), $M(g)$ est aussi vraie, donc l'assertion (ii) résulte de
(6.8.3, (i)). Quant à l'assertion (i), elle résulte directement de (6.6.1, (i)).
\end{proof}

\begin{theorem}[6.8.6]
\label{IV.2.s6.8.6-fr}
Soient $X$, $Y$ deux préschémas localement noethériens, $f:X\to Y$ un
morphisme localement de type fini, $x$ un point de $X$, $y=f(x)$. Les
propriétés suivantes sont équivalentes:
\begin{enumerate}
\item[a)] $f$ est lisse au point $x$.
\item[b)] $f$ est régulier au point $x$.
